\subsection{Inclusion of valuation rings of the same field}

PROPOSITION 1. Let K be a field and A a valuation ring of K. Then:

\begin{enumerate}
\def\labelenumi{(\alph{enumi})}
\item
  Every ring \(\mathbf{B}\) such that \(\mathbf{A} \subset \mathbf{B} \subset \mathbf{K}\) is a valuation ring of \(\mathbf{K}\) ;
\item
  The maximal ideal \(\mathfrak{m}(\mathbf{B})\) of such a ring is contained in \(\mathbf{A}\) and it is a prime ideal of \(\mathbf{A}\) ;
\item
  The mapping \(\mathfrak{p} \mapsto \mathbf{A}_{\mathfrak{p}}\) is a decreasing bijection of the set of prime ideals of \(\mathbf{A}\) onto the set of rings \(\mathbf{B}\) such that \(\mathbf{A} \subset \mathbf{B} \subset \mathbf{K}\) ; its inverse bijection is the mapping \(\mathbf{B} \mapsto \mathfrak{m}(\mathbf{B})\) .
\end{enumerate}

If B is a ring such that \(A \subset B \subset K\) and \(x \in K - B\) , then \(x \in K - A\) , whence \(x^{-1} \in \mathfrak{m}(A) \subset B\) , which proves both that B is a valuation ring of K and that \(\mathfrak{m}(B) \subset \mathfrak{m}(A)\) ; as \(\mathfrak{m}(B) = \mathfrak{m}(B) \cap A\) is a prime ideal of A, we have shown (a)

and (b). Moreover, \(\mathbf{A}_{\mathfrak{m}(\mathbf{B})} \subset \mathbf{B}\) ; conversely, if \(x \in \mathbf{B} - \mathbf{A}\) , then \(x^{-1} \in \mathbf{A}\) and \(x^{-1} \notin \mathfrak{m}(\mathbf{B})\) and hence \(x \in \mathbf{A}_{\mathfrak{m}(\mathbf{B})}\) ; thus \(\mathbf{A}_{\mathfrak{m}(\mathbf{B})} = \mathbf{B}\) . Finally let \(\mathfrak{p}\) be a prime ideal of \(\mathbf{A}\) ; we write \(\mathbf{B} = \mathbf{A}_{\mathfrak{p}}\) ; then \(\mathfrak{m}(\mathbf{B}) \cap \mathbf{A} = \mathfrak{p}\) . (Chapter II, §2, no. 5, Proposition 11) and \(\mathfrak{m}(\mathbf{B}) \subset \mathbf{A}\) by (b); hence \(\mathfrak{m}(\mathbf{B}) = \mathfrak{p}\) , which shows that the mappings \(\mathfrak{p} \mapsto \mathbf{A}_{\mathfrak{p}}\) and \(\mathbf{B} \mapsto \mathfrak{m}(\mathbf{B})\) of the statement are inverse bijections.

COROLLARY. The set of subrings of K containing A is totally ordered by inclusion.

The set of prime ideals of A is totally ordered by inclusion (§ 1, no. 2, Theorem 1 (e)) and the mapping \(\mathfrak{p} \mapsto \mathrm{A}_{\mathfrak{p}}\) reverses the inclusion relations.

PROPOSITION 2. Let K be a field, B a valuation ring of K and \(h_B\) the place of K associated with B (with values in \(\kappa(B)\) ). Then the mapping \(A \mapsto h_B(A)\) defines a bijection of the set \(\mathfrak{A}\) of valuation rings of K contained in B onto the set \(\mathfrak{A}'\) of valuation rings of \(\kappa(B)\) .

If \(\mathbf{A} \in \mathfrak{A}\) , then \(h_{\mathrm{B}}(\mathbf{A}) \in \mathfrak{A}'\) : for if \(x' = h_{\mathrm{B}}(x)\) (where \(x \in \mathbf{B}\) ) is an element of \(\kappa(\mathbf{B}) - h_{\mathrm{B}}(\mathbf{A})\) , then \(x \notin \mathbf{A}\) , hence \(x^{-1} \in \mathbf{A}\) and \(h_{\mathrm{B}}(x)^{-1} \in h_{\mathrm{B}}(\mathbf{A})\) . On the other hand, for \(\mathbf{A} \in \mathfrak{A}\) , \(\mathbf{A} \supset \mathfrak{m}(\mathbf{B})\) (Proposition 1 (b)) and hence the mapping, \(\mathbf{A} \mapsto h_{\mathrm{B}}(\mathbf{A})\) is injective. Finally, let \(\mathbf{A}' \in \mathfrak{A}'\) and \(\mathbf{A} = \bar{h}_{\mathrm{B}}^{-1}(\mathbf{A}') \subset \mathbf{B}\) ; we shall show, which will complete the proof, that \(\mathbf{A} \in \mathfrak{A}\) ; if \(x \in \mathbf{K} - \mathbf{A}\) , then either \(x \notin \mathbf{B}\) , or \(x \in \mathbf{B}\) ; if \(x \notin \mathbf{B}\) , then \(x^{-1} \in \mathfrak{m}(\mathbf{B}) \subset \mathbf{A}\) ; if \(x \in \mathbf{B}\) , then \(h_{\mathrm{B}}(x) \in \kappa(\mathbf{B})\) and \(h_{\mathrm{B}}(x) \notin \mathbf{A}'\) , hence \(h_{\mathrm{B}}(x^{-1}) \in \mathbf{A}'\) and we conclude again that \(x^{-1} \in \mathbf{A}\) ; hence \(\mathbf{A} \in \mathfrak{A}\) .

COROLLARY. Let A and B be two valuation rings of K, where \(A \subset B\) ; let \(A' = h_{B}(A)\) , which is a valuation ring of \(\kappa(B)\) . The residue field \(\kappa(A')\) of \(A'\) is canonically isomorphic to the residue field \(\kappa(A)\) of A and the place \(h_{A}\) associated with A is the composition \(h_{A'} \circ h_{B}\) of the places associated with \(A'\) and B.

Since the local ring A' is a quotient of the local ring A, their residue fields are canonically isomorphic and the equation \(h_{\mathrm{A}}(x) = h_{\mathrm{A}'}(h_{\mathrm{B}}(x))\) holds for \(x \in A\) . On the other hand, if \(x \in B - A\) , then \(h_{\mathrm{B}}(x) \notin A'\) and the two sides of the equation are equal to \(\infty\) ; the same is true if \(x \in K - B\) .

Remark. Conversely, let \(f\) be a place of \(\mathbf{K}\) with values in \(\mathbf{K}'\) and \(f'\) a place of \(\mathbf{K}'\) with values in \(\mathbf{K}''\) . Then \(f' \circ f\) is a place of \(\mathbf{K}\) whose ring is contained in the ring of the place \(f\) .

